\section{The third homology of the real attachment}\label{sec:third-homology}

\subsection{The real attachment and its chains}

We continue with $W$ and the two marked core coverings of
Section~\ref{sec:second-homology}. Put $(d_0,d_1)=(3,4)$ and use the
unimodular coordinate matrices $\mathcal B_i$ of \eqref{mt:eq:B}.
Explicitly, in the common $e$ marking they are
\begin{equation}\label{th:eq:B}
\mathcal B_0=\begin{pmatrix}0&-1&-1&4\\0&0&-1&2\\1&1&1&-6\\0&0&0&1\end{pmatrix},
\qquad
\mathcal B_1=\begin{pmatrix}1&1&0&-1\\0&1&0&-1\\0&-1&1&3\\0&0&0&1\end{pmatrix}.
\end{equation}
Write $(y_i,t)^t=\mathcal B_i^{-1}a$, so $t=a_4$.

All mapping-torus fundamental classes have interval-first orientation.
The boundary maps and their core restrictions are
\begin{equation}\label{th:eq:maps}
\begin{aligned}
j_0[a,v]&=([y_0,3t+v],e^{-2\pi it}),&
f_0(a)&=[y_0,3t],\\
j_1[a,v]&=([y_1,4t-v],e^{2\pi it}),&
f_1(a)&=[y_1,4t].
\end{aligned}
\end{equation}

The double mapping cylinder \eqref{mt:eq:double} is homotopy
equivalent to $W$ by the preceding collar retractions. Its meridian
homotopies remain $[y_0,3t+v]$ and $[y_1,4t-v]$. In particular,
the positive cusp path first traverses $\gamma_0^{-1}$ and then
$\gamma_1^{-1}$, with the joining arc between them.

These maps admit finite compatible triangulations.
Intersect the finitely many $P_i$-translates of a rational torus subdivision.
Triangulate that intersection and use its prisms for $G_i$.
Cut the source of $f_i$ at integral values of $d_it$ and at inverse images of target faces.
Common refinements give cellular maps and the double mapping cylinder.

Every additional rational subtorus below can be included by the same refinement.
There is no requirement for one finite triangulation invariant under both $T_0$ and $T_1$.

We use the chain cone
\begin{equation}\label{th:eq:cone}
\begin{aligned}
\mathcal D_n&=C_n(G_0)\oplus C_n(G_1)\oplus C_{n-1}(F),\\
d(a_0,a_1,b)&=(da_0+f_{0*}b,\,da_1-f_{1*}b,\,-db).
\end{aligned}
\end{equation}
Here $C_*$ can be the compatible cellular chains or piecewise smooth singular chains.
Subdivision and prism comparisons identify the two diagrams.
The shifted term uses the negative of the interval-first cylinder orientation.
The resulting exact sequence contains
\begin{equation}\label{th:eq:sequence}
0\longrightarrow\operatorname{coker}J_3
\longrightarrow H_3(W)\xrightarrow{\partial}\ker J_2
\longrightarrow0,\qquad J_n=(f_{0*},-f_{1*}).
\end{equation}
The connecting map sends a cone cycle to $[b]$.

\subsection{Integral three-cycles in the finite cores}

Write $x_1,x_2,x_3$ for the separate ordered coordinate bases of $X_i$.
Let $A_i\in H_3(G_i)$ be the oriented fiber three-torus.
The invariant oriented planes have primitive bivectors
\[
w_0=x_1x_2-x_1x_3+x_2x_3,\qquad w_1=x_1x_2.
\]
Juxtaposition denotes exterior product.

Their primitive sublattices are
$\langle x_1-x_3,x_2-x_3\rangle$ and $\langle x_1,x_2\rangle$.
The matrices $P_i$ preserve these sublattices with determinant one.
Let $S_i$ be the fundamental class of the positive core mapping torus of that plane.

The core Wang sequence is
\[
0\longrightarrow
\operatorname{coker}(\bigwedge^3P_i-I)
\longrightarrow H_3(G_i)
\longrightarrow\ker(\bigwedge^2P_i-I)
\longrightarrow0.
\]
It follows directly from the interval-prism differential.
Since $\det P_i=1$, its left group is $\mathbb Z A_i$.
Solving the exterior-square equations gives $\mathbb Zw_i$ on the right.
The class $S_i$ maps to $w_i$.
Thus $(A_i,S_i)$ is an integral basis of $H_3(G_i)$.

In the adapted source basis $(x_1,x_2,x_3,t)$, the covering matrices are
\[
F^{\mathrm{ad}}_{0,3}=
\begin{pmatrix}1&0&0&0\\0&1&-1&1\end{pmatrix},\qquad
F^{\mathrm{ad}}_{1,3}=
\begin{pmatrix}1&0&0&0\\0&4&2&-2\end{pmatrix}.
\]
Their columns are the lexicographic exterior triples.
To verify the second rows, divide the covering interval into $d_i$ pieces.
On a mixed three-torus, the pieces contribute the norm
$\sum_{r=0}^{d_i-1}\bigwedge^2P_i^r$.
Moving the last interval to the first position contributes sign $+1$.
The two norm matrices are the columns $w_i$ times the rows
$(1,-1,1)$ and $(4,2,-2)$, respectively.

The first row on a mixed three-torus is zero:
the invariant volume form $dy_1\wedge dy_2\wedge dy_3$
has only two independent fiber derivatives there.
It integrates to one on $A_i$ and zero on $S_i$.
This excludes an additional fiber-volume coefficient.

Transforming by $\bigwedge^3\mathcal B_i^{-1}$ gives the actual fiber maps
\begin{equation}\label{th:eq:F3}
F_{0,3}=\begin{pmatrix}1&6&2&-4\\0&3&1&-2\end{pmatrix},\qquad
F_{1,3}=\begin{pmatrix}1&-3&-1&1\\0&6&2&-4\end{pmatrix}.
\end{equation}
Consequently
\begin{equation}\label{th:eq:J3}
J_3=\begin{pmatrix}
1&6&2&-4\\0&3&1&-2\\-1&3&1&-1\\0&-6&-2&4
\end{pmatrix}
\end{equation}
in target basis $(A_0,S_0,A_1,S_1)$.
Its second column is three times its third.
Its first, third, and fourth columns, followed by $(0,0,0,1)^t$,
form a determinant-one matrix.
Hence
\begin{equation}\label{th:eq:c}
\operatorname{coker}J_3=\mathbb Zc,\qquad
c=[S_1],\qquad q=(0,2,0,1),\qquad q(c)=1.
\end{equation}
The primitive row $q$ annihilates $J_3$.

We also need the degree-two fiber relations.
The degree-two core bases consist of one fiber torus and the mapping torus
of $U_0=(1,1,1)^t$ or $U_1=(-1,-1,2)^t$.
The fiber quotient rows are $(1,-1,1)$ and $(2,1,-1)$.
The mixed-torus coefficients are the negative norms of $P_i$ on $H_1(X_i)$.
Invariant vertical two-forms exclude a residual fiber coefficient, as above.

This gives
\begin{equation}\label{th:eq:J2}
J_2=
\begin{pmatrix}
3&1&0&-2&0&0\\0&0&0&0&0&-1\\
-3&-1&0&2&-3&-1\\0&0&0&0&2&2
\end{pmatrix}.
\end{equation}
Its source basis is $(e_1e_2,e_1e_3,e_1e_4,e_2e_3,e_2e_4,e_3e_4)$.
An integral vector $y$ belongs to its kernel exactly when
\begin{equation}\label{th:eq:kerJ2}
y_5=y_6=0,\qquad 3y_1+y_2-2y_4=0.
\end{equation}

\subsection{Four actual meridian three-manifolds}

Use the ordered fiber basis
\[
\alpha_1=e_1-e_3,\quad \alpha_2=e_3-e_2,\quad
\beta_1=e_3,\quad\beta_2=e_4.
\]
It is unimodular, and
$T_\infty\alpha_j=\alpha_j$, $T_\infty\beta_j=\beta_j+\alpha_j$.
Consider the four primitive invariant subtori with oriented bivectors
\begin{equation}\label{th:eq:v}
v_1=\alpha_1\alpha_2,\quad
v_2=\alpha_1\beta_1,\quad
v_3=\alpha_2\beta_2,\quad
v_4=(\alpha_1+\alpha_2)(\beta_1+\beta_2).
\end{equation}

Let $V_j$ be the corresponding oriented torus and $M_j=[M(V_j)]$
its positive mapping-torus fundamental class in $H_3(B_{T_\infty})$.
The first restricted monodromy is the identity.
Each of the other three has matrix
$\left(\begin{smallmatrix}1&1\\0&1\end{smallmatrix}\right)$
in the displayed ordered basis.
Thus these are actual closed oriented three-manifolds.

The following chain construction retains their nontrivial seams.
Choose a finite oriented fundamental surface chain $u_j$ in $V_j$.
There is a piecewise smooth singular three-chain $a_j$, supported in $V_j$, with
\begin{equation}\label{th:eq:seam}
da_j=T_{\infty*}u_j-u_j.
\end{equation}

To construct it, take a common refinement of the triangulations for $u_j$
and $T_{\infty*}u_j$.
Let $p_j,p'_j$ be their subdivision prisms to its same oriented fundamental chain $u'_j$.
Then $dp_j=u'_j-u_j$ and $dp'_j=u'_j-T_{\infty*}u_j$.

Set $a_j=p_j-p'_j$.
The prism maps stay in $V_j$.
This construction uses surface-chain subdivision, without a homotopy between
the identity map and the restricted unipotent automorphism.

Put $u_j^-=T_{0*}^{-1}u_j$.
Let $R_i$ be the first three rows of $\mathcal B_i^{-1}$.
The two inverse-meridian prisms, extended linearly over surface chains, are
\begin{equation}\label{th:eq:prisms}
Q_0(v,a)=[R_0a,3a_4-v],\qquad
Q_1(v,a)=[R_1a,4a_4+v].
\end{equation}
They define actual smooth maps $[0,1]\times F\to G_i$.
The period relations make these formulas independent of the lift of $a$.
Their endpoint identities give
\[
dQ_0(u_j)=f_{0*}(u_j^--u_j),\qquad
dQ_1(u_j^-)=f_{1*}(T_{\infty*}u_j-u_j^-).
\]

Thus the actual global cycle is
\begin{equation}\label{th:eq:globalcycle}
Z_j=\bigl(Q_0(u_j),\,Q_1(u_j^-)-f_{1*}a_j,\,u_j-u_j^-\bigr)
\in\mathcal D_3.
\end{equation}
Substituting \eqref{th:eq:seam} into \eqref{th:eq:cone} gives $dZ_j=0$
in each component, with the displayed signs.

Cutting $M(V_j)$ along its marked fiber gives its interval prism and the seam correction $-a_j$.
The spine retraction takes these pieces to \eqref{th:eq:globalcycle}.
Therefore $[Z_j]=b_{\infty*}M_j$ for the actual boundary inclusion.

The connecting vectors are
$\partial[Z_j]=(I-\bigwedge^2T_0^{-1})v_j$.
Their complete matrix is
\begin{equation}\label{th:eq:D}
D_2=
\begin{pmatrix}
-1&0&0&-1\\3&2&0&3\\0&0&2&3\\
0&1&0&0\\0&0&0&0\\0&0&0&0
\end{pmatrix}.
\end{equation}
Its columns $d_j$ satisfy \eqref{th:eq:kerJ2}.
The three vectors $d_1,d_2,d_3-d_4$ form an integral basis of that kernel.
Indeed, in those coordinates an element is
\[
(-r+t,\,3r+2s-3t,\,-t,\,s,\,0,\,0),
\qquad (r,s,t)\in\mathbb Z^3,
\]
and \eqref{th:eq:kerJ2} gives the unique inverse coordinates.

Solving $D_2n=0$ over $\mathbb Z$ gives
\begin{equation}\label{th:eq:r}
\ker D_2=\mathbb Z(-2,0,-3,2)^t.
\end{equation}
The corresponding combination of the $Z_j$ lies in the core subgroup $\mathbb Zc$.
Its coefficient is the remaining integral question.

\subsection{A closed cochain determines the primitive relation}

On $X_0$ and $X_1$, respectively, put
\[
\omega_0=dy_1\wedge dy_2-dy_1\wedge dy_3+dy_2\wedge dy_3,\qquad
\omega_1=2dy_1\wedge dy_2+dy_1\wedge dy_3-dy_2\wedge dy_3.
\]
Each is $P_i$-invariant.
Define the closed core forms
\begin{equation}\label{th:eq:theta}
\theta_0=\tfrac23\,d\tau\wedge\omega_0,\qquad
\theta_1=\tfrac12\,d\tau\wedge\omega_1.
\end{equation}
The integrals of $\omega_0,\omega_1$ on the invariant planes $w_0,w_1$
are $3$ and $2$.
Consequently the pair of forms evaluates on
$(A_0,S_0,A_1,S_1)$ as $(0,2,0,1)$.
Its restriction to the core quotient is the primitive homomorphism $q$.

The actual pullbacks agree as differential forms on $F$:
\begin{equation}\label{th:eq:eta}
f_0^*\theta_0=f_1^*\theta_1=\eta,\qquad
\eta=6\,da_1\wedge da_2\wedge da_4
 +2\,da_1\wedge da_3\wedge da_4
 -4\,da_2\wedge da_3\wedge da_4.
\end{equation}
Here $a_1,\ldots,a_4$ are the real coordinates in the $e$ basis.
This equality follows by substituting \eqref{th:eq:B} into \eqref{th:eq:theta}.
It is an equality of constant forms, including all coefficients.

On the piecewise smooth chain cone define
\[
\mathcal L(a_0,a_1,b)=\int_{a_0}\theta_0+\int_{a_1}\theta_1.
\]
For a boundary, Stokes' formula gives zero on the two core differentials.
The remaining terms are
$\int_b f_0^*\theta_0-\int_b f_1^*\theta_1=0$.
Thus $\mathcal L$ is a real homology functional.
It need not be integral on every global class.
Its restriction to the integral core subgroup is exactly $q$.

Each $a_j$ in \eqref{th:eq:seam} is supported in a two-dimensional torus.
Therefore $\int_{a_j}\eta=0$.

In the first prism, the meridian derivative of the core parameter is $-1$.
In the second prism it is $+1$.
Their fiber derivatives give
\begin{equation}\label{th:eq:evaluation}
\mathcal L(Z_j)
=-\tfrac23\,\omega_0\bigl((\bigwedge^2R_0)v_j\bigr)
 {}+\tfrac12\,\omega_1\bigl((\bigwedge^2(R_1T_0^{-1}))v_j\bigr).
\end{equation}
In the $e$ exterior basis, this evaluation row is
\[
\lambda=(-\tfrac12,-\tfrac16,-\tfrac32,\tfrac13,0,-1).
\]
Its four values on \eqref{th:eq:v} are
\begin{equation}\label{th:eq:values}
(\mathcal L(Z_1),\mathcal L(Z_2),\mathcal L(Z_3),\mathcal L(Z_4))
=(0,-\tfrac16,-1,-2).
\end{equation}
The fractional second value is compatible with the stated real coefficient field.

Using \eqref{th:eq:r}, the integral core class
$-2[Z_1]-3[Z_3]+2[Z_4]$ has evaluation $-1$.
Since $q(c)=1$ on the free core subgroup, this proves the integral relation
\begin{equation}\label{th:eq:relation}
-2[Z_1]-3[Z_3]+2[Z_4]=-c
\quad\hbox{in }H_3(W;\mathbb Z).
\end{equation}
In particular, the relation produces a primitive core class.
No intersection or linking pairing is needed for this calculation.

\subsection{The integral attaching matrix}

Set
$t_1=[Z_1]$, $t_2=[Z_2]$, and $t_3=[Z_3]-[Z_4]$.
Their connecting classes form the integral basis of $\ker J_2$ given above.
They therefore split \eqref{th:eq:sequence}.
Together with $c$, they are an integral basis of $H_3(W)$.

Equation~\eqref{th:eq:relation} now gives
\begin{equation}\label{th:eq:meridianmatrix}
b_{\infty*}(M_1,M_2,M_3,M_4)=
\begin{pmatrix}
1&0&-2&-2\\0&1&0&0\\0&0&-2&-3\\0&0&1&1
\end{pmatrix}
\end{equation}
in the target basis $(t_1,t_2,t_3,c)$.
The determinant is one.
This proves both the image and kernel over the integers.

We state the marked cusp input for the complete attachment.
Let $N_c$ be a compact smooth cusp neighborhood with boundary $B=B_{T_\infty}$.
Let $k:B\to N_c$ be its actual marked inclusion.

Assume its degree-three homology is free of rank two, and that
\begin{equation}\label{th:eq:local}
H_3(B)=
\bigoplus_{j=1}^4\mathbb ZM_j\oplus\mathbb ZC_1\oplus\mathbb ZC_2,
\qquad
k_*M_j=0,\quad k_*C_i=n_i,
\end{equation}
where $(n_1,n_2)$ is an integral basis of $H_3(N_c)$.
The fiber classes $C_1,C_2$ are the actual oriented three-tori
$\alpha_1\beta_1\beta_2$ and $\alpha_2\beta_1\beta_2$.
These are the local data of the marked toric specialization with periods
$(qc_1,c_2),(c_3,qc_4)$ and its retained section.
Theorem~\ref{ga:thm:boundary} supplies them, and
\eqref{ga:eq:comparison} preserves the actual meridian homotopy.

For clarity, the boundary-group splitting in \eqref{th:eq:local}
has a direct integral check.
The four $v_j$ form a basis of $\ker(\bigwedge^2T_\infty-I)$.
The fiber quotient in degree three is freely generated by $C_1,C_2$.
The actual $M_j$ project to $v_j$ in the Wang sequence.
For the marked toric collapse, their images have two-dimensional carriers,
and the fiber map sends $C_i$ to $n_i$.
These geometric properties are the local input in \eqref{th:eq:local}.

The fiber inclusion into $W$ lands in the core quotient.
The evaluation \eqref{th:eq:eta} gives
$b_{\infty*}C_1=2c$ and $b_{\infty*}C_2=4c$.
Consequently the full degree-three attaching map is
\begin{equation}\label{th:eq:full}
(b_{\infty*},-k_*)_3=
\begin{pmatrix}
1&0&-2&-2&0&0\\
0&1&0&0&0&0\\
0&0&-2&-3&0&0\\
0&0&1&1&2&4\\
0&0&0&0&-1&0\\
0&0&0&0&0&-1
\end{pmatrix}.
\end{equation}
The source basis is $(M_1,M_2,M_3,M_4,C_1,C_2)$.
The target basis is $(t_1,t_2,t_3,c,n_1,n_2)$.

\begin{theorem}\label{th:thm:third}
For the specified real fillings,
$H_3(W;\mathbb Z)=\mathbb Z^4$, with basis $(t_1,t_2,t_3,c)$.
The actual meridian map \eqref{th:eq:meridianmatrix} is an integral isomorphism.
Under the marked cusp data \eqref{th:eq:local}, the full map
\eqref{th:eq:full} is an integral isomorphism.
Both its kernel and cokernel vanish.
\end{theorem}

\begin{proof}
The global sequence, the basis of its quotient, and the primitive relation
give the stated integral basis.
The meridian determinant is one.
The last two diagonal entries of the full block matrix are both $-1$,
so its determinant is also one.
\end{proof}

To determine $H_3(Y)$ for $Y=W\cup_BN_c$, the preceding attaching map
in degree two must also be injective.
Theorem~\ref{ga:thm:boundary} gives a surjective $k_{*,2}$ whose
kernel has basis the two angular meridian tori $z_i=M(\alpha_i)$.
Their images in the finite union can be checked integrally as follows.
The one-dimensional version of the actual prisms \eqref{th:eq:prisms}
gives their connecting images $(I-T_0^{-1})\alpha_i=e_1,2e_1$.

The difference $z_2-2z_1$ has fiber circle
$\delta=-2e_1-e_2+3e_3$, fixed by both finite monodromies.
Its core images are the negative $U_0$ mapping torus
and the positive $U_1$ mapping torus.
In the degree-two core basis, their sum is $c_2=(0,-1,0,1)^t$.
Columns $2,5,6$ of \eqref{th:eq:J2}, followed by $c_2$, have determinant one.

The degree-one fiber map, obtained directly from \eqref{th:eq:maps}, is
\[
J_1=\begin{pmatrix}0&0&1&6\\0&0&0&3\\0&-1&-1&2\\0&0&0&-4\end{pmatrix}.
\]
Its rows give the integral kernel $\mathbb Ze_1$.
Thus the two actual angular meridian tori map isomorphically onto $H_2(W)$.

If an element of $H_2(B)$ is killed by both $b_{\infty*}$ and $k_*$,
it belongs to their two-dimensional meridian subgroup and is zero.
This proves the needed injectivity of $(b_{\infty*},-k_*)_2$.

\begin{corollary}\label{th:cor:Y}
For the specified marked toric attachment,
\[
H_3(Y;\mathbb Z)=0.
\]
In particular, this group has no torsion.
\end{corollary}

\begin{proof}
The Mayer--Vietoris sequence gives
\[
\operatorname{coker}(b_{\infty*},-k_*)_3
\longrightarrow H_3(Y)
\longrightarrow\ker(b_{\infty*},-k_*)_2.
\]
Both outer groups vanish by the actual matrices and the preceding argument.
\end{proof}

\subsection{The degree-four meridian map}

The degree-three isomorphism determines the middle homology, but a
chain model of the entire attachment must also retain the fundamental
four-classes of the finite cores and the cusp. Their signs depend on
the interval-first orientation already used in \eqref{th:eq:maps}.
Let $D_i\in H_4(G_i;\mathbb Z)$ be the corresponding fundamental class,
and let $F_e=e_1e_2e_3e_4$ be the oriented fiber class.
The coordinate changes $\mathcal B_i$ have determinant one.
Moving the last coordinate $t$ to the first position changes orientation
by $-1$, so the two coverings send $F_e$ to $-3D_0$ and $-4D_1$.
Consequently
\begin{equation}\label{th:eq:J4}
 J_4=\binom{-3}{4},\qquad
 c_4=D_0-D_1,\qquad q_4=(4,3).
\end{equation}
The columns $J_4,c_4$ have determinant $-1$; moreover
$q_4J_4=0$ and $q_4(c_4)=1$.
Thus $\operatorname{coker}J_4=\mathbb Zc_4$, with no torsion.
The row equations of \eqref{th:eq:J3} give
\[
 \ker J_3=\mathbb Zk_3,\qquad k_3=e_1e_2e_4-3e_1e_3e_4.
\]
Indeed they force the coefficients of $e_1e_2e_3$ and $e_2e_3e_4$
to vanish, and the coefficient of $e_1e_3e_4$ to be minus three times
that of $e_1e_2e_4$. The cone sequence therefore gives
\begin{equation}\label{th:eq:H4-extension}
 0\longrightarrow\mathbb Zc_4\longrightarrow H_4(W;\mathbb Z)
 \xrightarrow{\partial}\mathbb Zk_3\longrightarrow0,
 \qquad H_5(W;\mathbb Z)=\ker J_4=0.
\end{equation}

The two primitive invariant three-vectors
\[
 w_1^{(3)}=\alpha_1\alpha_2\beta_1,\qquad
 w_2^{(3)}=\alpha_1\alpha_2\beta_2
\]
define subtori $V_j^{(3)}\subset F$. Their sublattices are generated by
subsets of the unimodular $(\alpha,\beta)$ basis, and the restricted
transport has determinant one. Put $L_j=[M(V_j^{(3)})]\in H_4(B_{T_\infty})$,
using the fundamental classes of these actual closed four-manifolds,
with interval-first orientation.
Choose fundamental three-chains $u_j^{(3)}$ and subdivision seam chains
$a_j^{(4)}$ supported on those three-tori, with
$da_j^{(4)}=T_{\infty*}u_j^{(3)}-u_j^{(3)}$.
The construction \eqref{th:eq:seam}--\eqref{th:eq:globalcycle}, now in
degree four, gives their images under the same actual inclusion
$b_\infty$. Its connecting map is still $I-T_{0*}^{-1}$ on the fiber
chain. In the lexicographic $e$-basis of $\bigwedge^3\Lambda$,
\begin{equation}\label{th:eq:D3}
 \begin{pmatrix}w_1^{(3)}&w_2^{(3)}\end{pmatrix}
 =\begin{pmatrix}-1&0\\0&-1\\0&1\\0&-1\end{pmatrix},
 \qquad
 (I-\bigwedge^3T_0^{-1})
 \begin{pmatrix}w_1^{(3)}&w_2^{(3)}\end{pmatrix}
 =\begin{pmatrix}0&0\\0&-1\\0&3\\0&0\end{pmatrix}.
\end{equation}
Hence $\partial b_{\infty*}L_1=0$ and
$\partial b_{\infty*}L_2=-k_3$.

To determine the first class inside the integral core subgroup, put
$\operatorname{vol}_i=dy_1\wedge dy_2\wedge dy_3$ on $X_i$.
Since $\det P_i=1$, the four-forms
\[
 \Theta_0=4\,d\tau\wedge\operatorname{vol}_0,\qquad
 \Theta_1=3\,d\tau\wedge\operatorname{vol}_1
\]
descend to the cores and are closed. They evaluate on $(D_0,D_1)$ as
$q_4$ and have identical pullbacks
$-12\,da_1\wedge da_2\wedge da_3\wedge da_4$ under $f_0,f_1$.
Their sum of integrals on the two core terms is therefore a closed
functional on \eqref{th:eq:cone}, by the same Stokes calculation as
\eqref{th:eq:eta}. Its restriction to $\mathbb Zc_4$ is $q_4$.
The two inverse-meridian prisms \eqref{th:eq:prisms} give the row
\begin{equation}\label{th:eq:lambda4}
 -4\bigwedge^3R_0+3\bigwedge^3(R_1T_0^{-1})
       =(-1,-33,-8,22).
\end{equation}
Its values on the columns of \eqref{th:eq:D3} are $1,3$.
The seam integrals vanish because their maps have three-dimensional
carriers. Since the first connecting image is zero and $q_4(c_4)=1$,
these evaluations prove
\[
 b_{\infty*}L_1=c_4,\qquad t_4:=b_{\infty*}L_2,\qquad
 H_4(W;\mathbb Z)=\mathbb Zc_4\oplus\mathbb Zt_4.
\]
The second class splits \eqref{th:eq:H4-extension} integrally because
its connecting image is the primitive generator $-k_3$.

The angular/deck fiber class is
$F_\alpha=\alpha_1\alpha_2\beta_1\beta_2=-F_e$.
The preceding common pullback consequently gives
$b_{\infty*}F_\alpha=12c_4$.
On the cusp side put $n_4=p_*F_\alpha$; it is the primitive top cell
$h\otimes a_{12}$ of \eqref{c26:eq:homologyp}.
Theorem~\ref{ga:thm:boundary} gives $k_*L_j=0$ and $k_*F_\alpha=n_4$.
Thus the full degree-four attachment is
\begin{equation}\label{th:eq:A4}
 \mathsf A_4:=(b_{\infty*},-k_*)_4
 =\begin{pmatrix}1&0&12\\0&1&0\\0&0&-1\end{pmatrix}
\end{equation}
in source basis $(L_1,L_2,F_\alpha)$ and target basis $(c_4,t_4,n_4)$.
Its determinant is $-1$.

\subsection{The complete integral cone}

Write $Y_{\mathrm{top}}=W\cup_BN_c$ for the actual marked union, where
$B=B_{T_\infty}$ and both inclusions are the specified boundary maps.
The established freeness of the piece homology now permits a reduction
of the full attachment. The chain contractions in
Section~\ref{sec:cusp-attachment} transport its attaching maps; the
following lemma gives the required compatibility for the global union.

\begin{lemma}\label{th:lem:integral-splitting}
Let $f:B_\bullet\to A_\bullet$ be a map of finite free integral chain
complexes whose homology groups are free. Then $\operatorname{Cone}(f)$
is chain-homotopy equivalent over $\mathbb Z$ to the cone of the induced
map $f_*:H_*(B)\to H_*(A)$, with zero differentials on these two
homology complexes. The equivalence transports the original map $f$.
\end{lemma}

\begin{proof}
Split the cycles and boundaries as in
Section~\ref{sec:cusp-attachment}: the two exact sequences with quotients
$\operatorname{im}d_{n}$ and $H_n$ split because these groups are free.
Thus each complex is its homology plus pairs on which the differential
is the identity. Let $i,p,h$ denote the resulting inclusion, projection
and contraction, so $pi=I$ and $I-ip=dh+hd$.
For $f_H=p_Afi_B$ and $U=h_Af+i_Ap_Afh_B$, direct substitution gives
\[
 f-i_Af_Hp_B=dU+Ud.
\]
With the cone convention of \eqref{th:eq:cone}, the triangular map
$(a,b)\mapsto(a+Ub,b)$ is an isomorphism from the cone of $f$ to the
cone of $i_Af_Hp_B$; its inverse subtracts $Ub$.
In the latter cone the identity pairs of $A$ and the shifted identity
pairs of $B$ split off. Contracting them leaves precisely the cone of
$f_H=f_*$. All these operations are integral.
\end{proof}

For this union the piece homology ranks, in degrees zero through five,
are
\begin{equation}\label{th:eq:piece-ranks}
\begin{array}{c|rrrrrr}
 &0&1&2&3&4&5\\\hline
 W&1&1&2&4&2&0\\
 N_c&1&2&4&2&1&0\\
 B_{T_\infty}&1&3&6&6&3&1.
\end{array}
\end{equation}
Every group in the table is free.
For $W$, this follows from Proposition~\ref{cc:prop:pi1},
Theorems~\ref{mt:thm:map} and~\ref{th:thm:third}, and
\eqref{th:eq:H4-extension}. The other two rows follow from
\eqref{c26:eq:homology} and Theorem~\ref{ga:thm:boundary}.
The groups in higher degrees vanish, since the finite models have
dimensions at most five, four and five, respectively.

Write $\mathsf A_n=(b_{\infty*},-k_*)_n$ for the induced attaching maps.
In the bases of \eqref{ga:eq:global1}, their first two matrices are
\[
 \mathsf A_0=\binom{1}{-1},\qquad
 \mathsf A_1=\begin{pmatrix}0&12&-1\\-1&0&0\\0&-1&0\end{pmatrix}.
\]
For a complete degree-two matrix, use the integral fiber lifts
\[
 \eta_0=\alpha_1\beta_2,\quad
 \eta_1=\alpha_1(\beta_1+\beta_2),\quad
 \eta_2=(\alpha_1+\alpha_2)\beta_2,\quad
 \eta_h=\beta_1\beta_2.
\]
Equation~\eqref{ga:eq:p23} sends these to $(E_0,E_1,E_2,h)$.
Put $s_2=b_{\infty*}z_1$ and retain the class
$c_2=(0,-1,0,1)$ used above. The row $(-2,-4,-2,-3)$ annihilates
\eqref{th:eq:J2} and evaluates to one on $c_2$.
The common fiber includes into $W$ through $f_0:F\to G_0$.
Its degree-two map is the first two rows of \eqref{th:eq:J2}.
Applying the first two entries of this quotient row to that map,
after expressing the four displayed lifts in the $e$-basis, gives
\[
 b_{\infty*}(\eta_0,\eta_1,\eta_2,\eta_h)=(-4,-6,0,4)c_2.
\]
Together with $b_{\infty*}z_2=2s_2+c_2$, this derives
\begin{equation}\label{th:eq:A2}
 \mathsf A_2=
 \begin{pmatrix}
 1&2&0&0&0&0\\
 0&1&-4&-6&0&4\\
 0&0&-1&0&0&0\\
 0&0&0&-1&0&0\\
 0&0&0&0&-1&0\\
 0&0&0&0&0&-1
 \end{pmatrix}
\end{equation}
in source basis $(z_1,z_2,\eta_0,\eta_1,\eta_2,\eta_h)$ and target
basis $(s_2,c_2,E_0,E_1,E_2,h)$.
The matrix $\mathsf A_3$ is \eqref{th:eq:full}, and $\mathsf A_4$ is
\eqref{th:eq:A4}. Their determinants, in increasing positive degree,
are $-1,1,1,-1$. The degree-five target is zero, so $\mathsf A_5=0$.

Apply Lemma~\ref{th:lem:integral-splitting} to the actual map in
\eqref{ga:eq:global}. The resulting finite integral complex has
\[
 \begin{aligned}
 \mathcal C_n&=H_n(W)\oplus H_n(N_c)\oplus H_{n-1}(B_{T_\infty}),\\
 d(a,b)&=(\mathsf A_{n-1}b,0),
 \end{aligned}
\]
where the first two summands are grouped together as $a$.
Its ranks are therefore
\begin{equation}\label{th:eq:complete-ranks}
 (\operatorname{rank}\mathcal C_0,\ldots,\operatorname{rank}\mathcal C_6)
 =(2,4,9,12,9,3,1).
\end{equation}
They are sums of the ranks in \eqref{th:eq:piece-ranks}.
The nonzero Smith factors of $d_1,\ldots,d_5$ are all one, with
multiplicities $1,3,6,6,3$; $d_6=0$.
Thus every middle group contracts in unit pairs. The remaining
degree-zero quotient is $\mathbb Z$, whose dual is the common constant
$1$ on both pieces. The remaining degree-six generator is the
suspended oriented boundary class; under the cone comparison it is
the orientation class of the glued manifold, with the boundary sign
fixed by the collar convention.

\begin{proposition}\label{th:prop:complete-cone}
The actual marked attachment has an integral chain complex
chain-homotopy equivalent to $\mathbb Z$ in degrees zero and six and
zero in other degrees. In particular
\[
 H_n(Y_{\mathrm{top}};\mathbb Z)=
 \begin{cases}\mathbb Z,&n=0,6,\\0,&1\leq n\leq5.\end{cases}
\]
After reduction modulo either $2$ or $3$, the differential ranks are
$(1,3,6,6,3,0)$ and the homology dimensions are $(1,0,0,0,0,0,1)$.
\end{proposition}

\begin{proof}
The preceding chain equivalence and unit eliminations give the integral
statement. Tensoring those equivalences with $\mathbb F_2$ or
$\mathbb F_3$ preserves each identity pair, giving the stated ranks.
The orientation generator and constant cohomology class are retained.
\end{proof}

The role of primitivity is visible in a different integral complex.
Replace the first column of \eqref{th:eq:A4} by twice that column,
leaving all other maps fixed. The new degree-four matrix is
\[
 \begin{pmatrix}2&0&12\\0&1&0\\0&0&-1\end{pmatrix}.
\]
It has a unit entry and a unit two-by-two minor, but determinant $-2$;
its Smith factors are $(1,1,2)$.
The resulting cone has $H_4=\mathbb Z/2$ and $H_5=0$.
Modulo $2$ its homology dimensions instead become
$(1,0,0,0,1,1,1)$: the degree-four class is
$H_4\otimes\mathbb F_2$, while the degree-five class is
$\operatorname{Tor}_1^{\mathbb Z}(H_4,\mathbb F_2)$.
Treating $2L_1$ as an integral basis element would therefore change
both the integral attachment and its coefficient reductions.
This complex is distinct from the cone of the geometric maps above.

\subsection{The smooth type of the specified real attachment}

The integral homology of the actual toric attachment is now that of a
six-sphere. To determine its smooth type one must also compute the
fundamental group. The first homology calculation alone does not suffice.

\begin{theorem}\label{th:thm:real-sphere}
Let $E\to\Sigma$ be the marked flat torus bundle above, with its two
specified finite fillings $N_0,N_1$.
Let $N_c$ be the toric filling of Section~\ref{sec:toric-cusp}. Glue them by their
displayed boundary diffeomorphisms. Then the resulting real manifold
\[
 Y_{\mathrm{top}}=E\cup_{j_0}N_0\cup_{j_1}N_1\cup_{j_\infty}N_c
\]
is diffeomorphic to the standard six-sphere.
\end{theorem}

\begin{proof}
The finite fillings are smooth disk bundles over compact torus mapping
tori. The toric filling is smooth because its maximal cones are
unimodular. Its tube is compact and its outer boundary is a smooth
torus bundle. The specified boundary maps reverse the boundary
orientations: a positive puncture meridian is opposite to the boundary
orientation of $\Sigma$. Gluing collars therefore gives a connected,
closed, oriented smooth six-manifold.

Write $W=E\cup N_0\cup N_1$ and
$K=\langle e_1-e_3,e_3-e_2\rangle$. By
Lemma~\ref{ga:lem:fundamental-group} and van Kampen,
$\pi_1(Y_{\mathrm{top}})$ is the group of
Proposition~\ref{cc:prop:pi1} modulo the normal closure of $K$ and
$\gamma_\infty$. This normal closure also contains
\[
 T_0(e_1-e_3)=e_1+e_2-e_3.
\]
The three vectors $e_1-e_3$, $e_3-e_2$, $e_1+e_2-e_3$ have
determinant $-1$ in $\mathbb Ze_1\oplus\mathbb Ze_2\oplus\mathbb Ze_3$.
Thus they kill $\ker\psi$ integrally.

If $c$ denotes the surviving
class of $e_4$, the conormal vectors in
\eqref{cc:eq:meridians} give
\[
 \gamma_0^3=c,\qquad \gamma_1^4=c^{-1}.
\]
Both transports fix $c$ in this quotient. Moreover,
$\gamma_\infty=\gamma_1^{-1}\gamma_0^{-1}$ is killed, so
$\gamma_1=\gamma_0^{-1}$. With $x=\gamma_0$ the group is
\[
 \langle c,x\mid[c,x]=1,\ x^3=c,\ x^4=c\rangle=1.
\]

Proposition~\ref{th:prop:complete-cone} gives
$H_j(Y_{\mathrm{top}};\mathbb Z)=0$ for $j=1,2,3$.
Poincar\'e duality and the universal coefficient theorem then give
$H_5=H^1=0$ and $H_4=H^2=0$, while $H_0=H_6=\mathbb Z$.
The manifold is therefore a simply connected integral homology
six-sphere. Milnor's smooth classification in dimensions five and six
\cite[\S9, Proposition B, pp.~109--110]{MilnorHC}
identifies it with the standard sphere.
\end{proof}

The finite maps above use $a_0=e_4/3$ and
$a_1=-(e_1+e_4)/4$ in the common basis. Their holomorphic
realization requires a comparison of the regular complex family
$\mathcal E_u\to\Sigma$ with $E$, preserving its section and all
three marked boundaries. Write this comparison as
$\Phi_u:\mathcal E_u\to E$. The local extensions must satisfy
\[
 \Psi_i j_i^{\mathrm{geom}}=j_i\Phi_u,
 \qquad \Psi_i:\mathcal N_{u,i}\to N_i,
 \qquad i=0,1,\infty,
\]
on whole collars.

The next two sections construct these maps for
the selected lifts and an adapted toric cusp. The finite comparison
retains the conormal factors $e=s^{-2}e'$ and $e=s^{-1}e'$.
The cusp comparison retains the full integral period lattice.

The broader question permits primitive increments
$a_0\in\frac13\Lambda^{T_0}$, $a_1\in\frac14\Lambda^{T_1}$
with $|12e_4^*(a_0+a_1)|=1$. Those choices require their own
marked attachments or a relative diffeomorphism to the ones used
here. The toric and complex-geometric constructions are described
in \cite[Sections 2--4]{Engel}.

